\chapter{Sprints 1 et 2~: Infrastructure et Gestion des Employés}

\noindent\textbf{Plan du chapitre}
\begin{enumerate}[nosep]
  \item Introduction
  \item Sprint 1~: Initialisation de l'infrastructure
  \item Sprint 2~: Gestion du cycle de vie des employés
  \item Conclusion
\end{enumerate}

% ─────────────────────────────────────────────────────────────
\section{Introduction}
% ─────────────────────────────────────────────────────────────

La première release établit les fondations sur lesquelles tout le reste repose.
Le sprint~1 met en place l'infrastructure complète~: serveur d'identité Keycloak,
base Oracle provisionnée par Flyway et squelettes des sept microservices. Le
sprint~2 livre le premier module métier~: la gestion complète du cycle de vie
des employés avec synchronisation bidirectionnelle Oracle--Keycloak. Ces deux
sprints constituent le prérequis indispensable à tout développement applicatif
sécurisé et déployable de façon reproductible.

% ─────────────────────────────────────────────────────────────
\section{Sprint 1~: Initialisation de l'infrastructure}
% ─────────────────────────────────────────────────────────────

\subsection{Objectifs du sprint}

\begin{table}[H]
\centering
\caption{Tâches et critères d'acceptation~: sprint 1}
\label{tab:s1_obj}
\begin{tabularx}{\textwidth}{|l|L|}
\hline
\rowcolor{headerblue!25}\textbf{Tâche} & \textbf{Critère d'acceptation} \\\hline
\rowcolor{rowgray}
Squelettes microservices & 7 projets Spring Boot compilables avec POM parent
  multimodule Maven~; Java 21 \\\hline
Keycloak 26 & Realm \texttt{gerai} opérationnel~; client \texttt{gerai-frontend}~;
  4 rôles realm~: \texttt{ADMIN\_RH}, \texttt{CHEF}, \texttt{EMPLOYE}, \texttt{COMITE} \\\hline
\rowcolor{rowgray}
Oracle 23ai + Flyway & Migrations V1--V4 appliquées~: tables \texttt{EMPLOYEES},
  \texttt{DEPARTMENTS}, \texttt{POSITIONS}, \texttt{LEAVE\_REQUESTS}~;
  vue \texttt{V\_ALL\_DEMANDES} \\\hline
Sécurité JWT & \texttt{JwtAuthenticationConverter} commun~: extrait
  \texttt{realm\_access.roles}, préfixe \texttt{ROLE\_} \\\hline
\rowcolor{rowgray}
Angular & Intégration \texttt{keycloak-angular}~; redirection login/logout~;
  garde de route par rôle \\\hline
Docker Compose & Démarrage en une commande~: Oracle (:1521), Keycloak (:8080),
  Kafka (:9092), Zookeeper (:2181), MailHog (:1025) \\\hline
\end{tabularx}
\end{table}

\subsection{Conception UML}

\subsubsection*{Cas d'utilisation~: S'authentifier}

\begin{figure}[H]
\centering
\begin{tikzpicture}

  % Acteur
  \draw[thick] (-2.5,4.2) circle(0.3cm);
  \draw[thick] (-2.5,3.9)--(-2.5,3.1) coordinate(wU);
  \draw[thick] (wU)++(0,.1)--++(-.38,-.33) (wU)++(0,.1)--++(.38,-.33);
  \draw[thick] (wU)--++(-.27,-.52) (wU)--++(.27,-.52);
  \node[font=\small] at (-2.5,2.2) {Utilisateur};

  % Frontière
  \begin{scope}[on background layer]
    \filldraw[draw=headerblue!60, fill=blue!3, rounded corners=5pt, dashed]
      (0,0.2) rectangle (9.5,5.6);
    \node[font=\bfseries\small, headerblue] at (4.75,5.3) {GerAI / Keycloak 26};
  \end{scope}

  % UCs
  \draw (3.8,4.6) ellipse(2.2cm and .48cm);
  \node[font=\scriptsize] at (3.8,4.6) {S'authentifier (PKCE)};

  \draw (7.5,4.6) ellipse(1.8cm and .48cm);
  \node[font=\scriptsize] at (7.5,4.6) {Obtenir le JWT};

  \draw (3.8,3.2) ellipse(2.4cm and .48cm);
  \node[font=\scriptsize] at (3.8,3.2) {Accéder au tableau de bord};

  \draw (3.8,1.8) ellipse(2.2cm and .48cm);
  \node[font=\scriptsize] at (3.8,1.8) {Se déconnecter};

  % Relations internes
  \draw[->,dashed,>=stealth] (5.95,4.6)--(5.3,4.6)
    node[midway,above,font=\scriptsize]{<<include>>};
  \draw[->,dashed,>=stealth] (3.8,4.12)--(3.8,3.68)
    node[midway,right,font=\scriptsize]{<<include>>};

  % Associations acteur
  \draw[->] (-2.18,4.2)--(1.6,4.6);
  \draw[->] (-2.18,4.0)--(1.6,3.2);
  \draw[->] (-2.18,3.8)--(1.6,1.8);

\end{tikzpicture}
\caption{Diagramme de cas d'utilisation raffiné~: S'authentifier}
\label{fig:uc_auth}
\end{figure}

\begin{table}[H]
\centering
\caption{Description textuelle du cas d'utilisation~: S'authentifier}
\label{tab:uc_auth_desc}
\begin{tabularx}{\textwidth}{|l|L|}
\hline
\rowcolor{headerblue!25}\textbf{Champ} & \textbf{Description} \\\hline
\rowcolor{rowgray}
Acteurs & Tout utilisateur~: \texttt{ADMIN\_RH}, \texttt{CHEF}, \texttt{EMPLOYE}, \texttt{COMITE} \\\hline
Précondition & L'utilisateur dispose d'un compte actif dans le realm \texttt{gerai} \\\hline
\rowcolor{rowgray}
Postcondition & JWT valide stocké en session~; rôle extrait et tableau de bord affiché \\\hline
Scénario nominal &
\begin{enumerate}[leftmargin=*,nosep]
  \item L'utilisateur navigue vers \texttt{/dashboard}
  \item Angular (\texttt{keycloak-angular}) détecte l'absence de session et redirige vers Keycloak (flux PKCE)
  \item L'utilisateur saisit ses identifiants sur la page Keycloak
  \item Keycloak retourne un authorization code~; Angular échange contre un \texttt{access\_token}
  \item Le JWT est injecté dans chaque requête HTTP via \texttt{HttpInterceptor}
  \item Le rôle détermine la route d'accueil~: \texttt{/admin}, \texttt{/chef} ou \texttt{/employe}
\end{enumerate} \\\hline
\rowcolor{rowgray}
Alternatives &
\begin{itemize}[leftmargin=*,nosep]
  \item Identifiants incorrects~: message Keycloak, nouvelle tentative possible
  \item Compte désactivé~: message « Compte inactif, contacter le RH »
  \item JWT expiré~: refresh automatique via \texttt{refresh\_token}
\end{itemize} \\\hline
\end{tabularx}
\end{table}

\subsubsection*{Diagramme de séquence~: Authentification OAuth2/OIDC avec PKCE}

\begin{figure}[H]
\centering
\begin{tikzpicture}[xscale=1.0, yscale=0.92]

  % Lifelines
  \foreach \x/\lbl in {0/Utilisateur, 3/Angular, 6/Keycloak 26, 9.5/Microservice}{
    \node[lifeline, minimum width=2.2cm] at (\x, 9) {\lbl};
    \draw[thick,dashed] (\x, 8.65)--(\x, 0);
  }

  \draw[seqmsg] (0.05,8.0) -- (2.9,8.0) node[midway,above]{\scriptsize Accès /dashboard};
  \draw[seqmsg] (3,7.2)    -- (5.9,7.2) node[midway,above]{\scriptsize redirect /auth?response\_type=code (PKCE)};
  \draw[seqret] (6,6.4)    -- (0.05,6.4)  node[midway,above]{\scriptsize Page de connexion Keycloak};
  \draw[seqmsg] (0.05,5.6) -- (5.9,5.6)  node[midway,above]{\scriptsize POST identifiants};
  \draw[seqret] (6,4.8)    -- (3.05,4.8) node[midway,above]{\scriptsize authorization\_code + redirect\_uri};
  \draw[seqmsg] (3,4.0)    -- (5.9,4.0)  node[midway,above]{\scriptsize POST /token (code + code\_verifier)};
  \draw[seqret] (6,3.2)    -- (3.05,3.2) node[midway,above]{\scriptsize access\_token + refresh\_token (JWT)};
  \draw[seqmsg] (3,2.4)    -- (9.45,2.4) node[midway,above]{\scriptsize GET /api/... Bearer \{jwt\}};
  \draw[seqmsg] (9.5,1.7)  -- (6.05,1.7) node[midway,above]{\scriptsize Vérification clé publique JWK};
  \draw[seqret] (6,1.0)    -- (9.45,1.0) node[midway,above]{\scriptsize Rôles extraits (ROLE\_EMPLOYE…)};
  \draw[seqret] (9.5,0.3)  -- (3.05,0.3) node[midway,above]{\scriptsize 200 OK + données};

\end{tikzpicture}
\caption{Diagramme de séquence~: flux OAuth2 PKCE avec Keycloak 26, realm \texttt{gerai}}
\label{fig:seq_auth}
\end{figure}

\subsection{Implémentation backend~: JwtAuthenticationConverter}

Chaque microservice partage la même configuration de sécurité OAuth2~:

\begin{lstlisting}[language=Java,caption={JwtAuthenticationConverter commun — SecurityConfig.java}]
@Bean
public JwtAuthenticationConverter jwtAuthenticationConverter() {
    JwtGrantedAuthoritiesConverter gc = new JwtGrantedAuthoritiesConverter();
    gc.setAuthoritiesClaimName("realm_access.roles"); // realm gerai
    gc.setAuthorityPrefix("ROLE_");
    JwtAuthenticationConverter conv = new JwtAuthenticationConverter();
    conv.setJwtGrantedAuthoritiesConverter(gc);
    return conv;
}
\end{lstlisting}

Les endpoints sont protégés par des annotations \texttt{@PreAuthorize}~:
\begin{lstlisting}[language=Java]
@GetMapping("/employes")
@PreAuthorize("hasAnyRole('ADMIN_RH', 'CHEF')")
public ResponseEntity<Page<EmployeeDTO>> list(Pageable p) { ... }
\end{lstlisting}

\subsection{Implémentation frontend~: initialisation Keycloak}

L'application Angular démarre après l'initialisation de Keycloak~:

\begin{lstlisting}[language=Java,caption={Initialisation Keycloak dans main.ts}]
bootstrapApplication(AppComponent, {
  providers: [
    provideKeycloak({
      config: { url: environment.keycloakUrl,
                realm: 'gerai', clientId: 'gerai-frontend' },
      initOptions: { onLoad: 'check-sso', silentCheckSsoRedirectUri:
        window.location.origin + '/assets/silent-check-sso.html' }
    })
  ]
});
\end{lstlisting}

\subsection{Résultats du sprint 1}

\begin{figure}[H]
  \centering
  % \includegraphics[width=0.75\textwidth]{images/sprint1_login.png}
  \caption{Interface de connexion Keycloak~: page générée par le realm \texttt{gerai}}
  \label{fig:s1_login}
\end{figure}

% ─────────────────────────────────────────────────────────────
\section{Sprint 2~: Gestion du cycle de vie des employés}
% ─────────────────────────────────────────────────────────────

\subsection{Objectifs du sprint}

\begin{table}[H]
\centering
\caption{Tâches et critères d'acceptation~: sprint 2}
\label{tab:s2_obj}
\begin{tabularx}{\textwidth}{|l|L|}
\hline
\rowcolor{headerblue!25}\textbf{Tâche} & \textbf{Critère d'acceptation} \\\hline
\rowcolor{rowgray}
Créer un employé & Code \texttt{EMP-NNNN} auto-incrémenté~; provisioning Keycloak~;
  e-mail avec identifiants provisoires envoyé via MailHog \\\hline
CRUD complet & Liste paginée avec filtres (département, statut, recherche
  plein texte via \texttt{searchByQuery()})~; édition~; désactivation \\\hline
\rowcolor{rowgray}
Sync bidirectionnelle & Désactivation Oracle → désactivation Keycloak (\texttt{enabled: false})
  en temps réel, sans intervention manuelle \\\hline
Profil enrichi & Compétences, langues, certifications, statistiques de congés
  (\texttt{soldeCongesRestants}) \\\hline
\rowcolor{rowgray}
Interface Angular & \texttt{ListeEmployesAdminComponent}, \texttt{FormEmployeAdminComponent},
  \texttt{ProfilEmployeAdminComponent} \\\hline
\end{tabularx}
\end{table}

\subsection{Conception UML}

\subsubsection*{Cas d'utilisation~: Gérer les employés}

\begin{figure}[H]
\centering
\begin{tikzpicture}

  % Acteur Admin
  \draw[thick] (-2.5,5.5) circle(0.3cm);
  \draw[thick] (-2.5,5.2)--(-2.5,4.4) coordinate(wA);
  \draw[thick] (wA)++(0,.1)--++(-.38,-.33) (wA)++(0,.1)--++(.38,-.33);
  \draw[thick] (wA)--++(-.27,-.52) (wA)--++(.27,-.52);
  \node[font=\small] at (-2.5,3.5) {Admin RH};

  % Frontière
  \begin{scope}[on background layer]
    \filldraw[draw=headerblue!60, fill=blue!3, rounded corners=5pt, dashed]
      (0,-0.2) rectangle (10.5,7.0);
    \node[font=\bfseries\small, headerblue] at (5.25,6.7) {employe-service~: gestion des employés};
  \end{scope}

  % UCs
  \draw (4,6.1) ellipse(2.3cm and .48cm);
  \node[font=\scriptsize] at (4,6.1) {Créer un employé};

  \draw (4,4.8) ellipse(2.8cm and .48cm);
  \node[font=\scriptsize,align=center] at (4,4.8) {Générer code EMP-NNNN (Oracle)};

  \draw (4,3.5) ellipse(2.8cm and .48cm);
  \node[font=\scriptsize,align=center] at (4,3.5) {Provisionner compte Keycloak};

  \draw (4,2.2) ellipse(2.8cm and .48cm);
  \node[font=\scriptsize,align=center] at (4,2.2) {Envoyer e-mail provisoire};

  \draw (4,0.9) ellipse(2.3cm and .48cm);
  \node[font=\scriptsize] at (4,0.9) {Désactiver un employé};

  \draw (8.2,4.8) ellipse(1.8cm and .48cm);
  \node[font=\scriptsize] at (8.2,4.8) {Valider le formulaire};

  % Relations internes
  \draw[->,dashed,>=stealth] (4,5.62)--(4,5.28) node[midway,right,font=\scriptsize]{<<include>>};
  \draw[->,dashed,>=stealth] (4,4.32)--(4,3.98) node[midway,right,font=\scriptsize]{<<include>>};
  \draw[->,dashed,>=stealth] (4,3.02)--(4,2.68) node[midway,right,font=\scriptsize]{<<include>>};
  \draw[->,dashed,>=stealth] (6.4,4.8)--(5.5,6.1) node[pos=0.4,above,font=\scriptsize]{<<extend>>};

  % Association acteur
  \draw[->] (-2.18,5.5)--(1.7,6.1);
  \draw[->] (-2.18,5.2)--(1.7,0.9);

\end{tikzpicture}
\caption{Diagramme de cas d'utilisation raffiné~: gérer les employés}
\label{fig:uc_emp}
\end{figure}

\begin{table}[H]
\centering
\caption{Description textuelle du cas d'utilisation~: Créer un employé}
\label{tab:uc_emp_desc}
\begin{tabularx}{\textwidth}{|l|L|}
\hline
\rowcolor{headerblue!25}\textbf{Champ} & \textbf{Description} \\\hline
\rowcolor{rowgray}Acteur & Administrateur RH (\texttt{ADMIN\_RH}) \\\hline
Précondition & L'administrateur est authentifié~; le département et le poste cibles existent \\\hline
\rowcolor{rowgray}Postcondition & Employé persité en Oracle~; compte Keycloak actif~; e-mail provisoire envoyé \\\hline
Scénario nominal &
\begin{enumerate}[leftmargin=*,nosep]
  \item L'admin remplit le formulaire~: nom, prénom, e-mail, département, poste, date d'embauche
  \item Angular valide les champs requis~; soumet \texttt{POST /api/employes}
  \item L'\textit{employe-service} génère le code \texttt{EMP-NNNN} (séquence Oracle)
  \item L'endpoint appelle \texttt{POST /admin/realms/gerai/users} (Keycloak Admin API)
  \item Le mot de passe provisoire est défini via \texttt{PUT …/reset-password}
  \item Le rôle est affecté via \texttt{PUT …/role-mappings/realm}
  \item Un e-mail de bienvenue est envoyé via MailHog avec les identifiants
\end{enumerate} \\\hline
\rowcolor{rowgray}Alternatives &
\begin{itemize}[leftmargin=*,nosep]
  \item E-mail déjà utilisé~: contrainte UNIQUE Oracle → erreur 409
  \item Keycloak indisponible~: transaction rollback~; employé non créé en base
\end{itemize} \\\hline
\end{tabularx}
\end{table}

\subsubsection*{Diagramme de séquence~: Créer un employé}

\begin{figure}[H]
\centering
\begin{tikzpicture}[xscale=1.0, yscale=0.92]

  \foreach \x/\lbl in {0/Admin, 2.6/Angular, 5.5/employe-service, 8.8/Keycloak 26, 12/MailHog}{
    \node[lifeline, minimum width=2.2cm] at (\x, 9) {\lbl};
    \draw[thick,dashed] (\x, 8.65)--(\x, 0);
  }

  \draw[seqmsg] (0.05,8.0)--(2.5,8.0) node[midway,above]{\scriptsize Remplissage formulaire};
  \draw[seqmsg] (2.6,7.2)--(5.4,7.2)  node[midway,above]{\scriptsize POST /api/employes \{firstName, email, deptId, …\}};

  % Activation employe-service
  \fill[umlblue!60] (5.38,6.9) rectangle (5.62,0.8);

  \draw[seqmsg] (5.55,6.2)--(5.55,6.2) node[right,font=\scriptsize]{Génération EMP-NNNN (séquence Oracle)};
  \draw[seqmsg] (5.55,5.5)--(8.7,5.5) node[midway,above]{\scriptsize POST /admin/realms/gerai/users \{username, email\}};
  \draw[seqret] (8.8,4.8)--(5.65,4.8) node[midway,above]{\scriptsize 201 keycloakUserId (UUID)};
  \draw[seqmsg] (5.55,4.1)--(8.7,4.1) node[midway,above]{\scriptsize PUT …/users/\{id\}/reset-password \{temporary: true\}};
  \draw[seqret] (8.8,3.4)--(5.65,3.4) node[midway,above]{\scriptsize 204 No Content};
  \draw[seqmsg] (5.55,2.7)--(8.7,2.7) node[midway,above]{\scriptsize PUT …/users/\{id\}/role-mappings/realm \{EMPLOYE\}};
  \draw[seqret] (8.8,2.0)--(5.65,2.0) node[midway,above]{\scriptsize 204 No Content};
  \draw[seqmsg] (5.55,1.3)--(11.9,1.3) node[midway,above]{\scriptsize SMTP — login + mot de passe provisoire};
  \draw[seqret] (5.5,0.6)--(2.65,0.6)  node[midway,above]{\scriptsize 201 \{employeeId: 42, code: "EMP-0042"\}};

\end{tikzpicture}
\caption{Diagramme de séquence~: création d'un employé et provisioning Keycloak (realm \texttt{gerai})}
\label{fig:seq_emp}
\end{figure}

\subsection{Implémentation backend~: code employé et sync Keycloak}

\begin{lstlisting}[language=SQL,caption={Génération du code EMP-NNNN via séquence Oracle}]
-- Migration V2 : séquence dédiée aux codes employés
CREATE SEQUENCE emp_code_seq START WITH 1 INCREMENT BY 1 NOCACHE NOCYCLE;

-- Dans EmployeeService.java :
Long nextVal = em.createNativeQuery(
    "SELECT emp_code_seq.NEXTVAL FROM DUAL")
    .getSingleResult();
employee.setEmployeeCode(String.format("EMP-%04d", nextVal));
\end{lstlisting}

La désactivation synchronise Oracle et Keycloak atomiquement~: si l'appel Keycloak
échoue, la transaction Oracle est annulée pour garantir la cohérence.

\subsection{Résultats du sprint 2}

\begin{figure}[H]
  \centering
  % \includegraphics[width=\textwidth]{images/sprint2_employes.png}
  \caption{Interface de gestion des employés~: liste filtrée avec codes EMP-NNNN et statuts}
  \label{fig:s2_list}
\end{figure}

\begin{figure}[H]
  \centering
  % \includegraphics[width=0.80\textwidth]{images/sprint2_profil.png}
  \caption{Vue profil d'un employé~: compétences, certifications, solde de congés}
  \label{fig:s2_profil}
\end{figure}

% ─────────────────────────────────────────────────────────────
\section{Conclusion}
% ─────────────────────────────────────────────────────────────

La première release livre trois garanties fondamentales~: un JWT Keycloak validé
de façon uniforme par les sept services via un \texttt{JwtAuthenticationConverter}
partagé~; un schéma Oracle reproductible en onze migrations Flyway~; une
synchronisation Oracle--Keycloak bidirectionnelle validée en sens création
(provisioning) et en sens désactivation (révocation immédiate). Ces garanties
conditionne l'ensemble des développements de la release suivante.
